% Fichier généré automatiquement par tools/generate_corrections.py.
% Source : STAT/STAT-03-Demarche/14_Sympact/14_Sympact.tex
% Statut : complet (3/3 question(s) renseignée(s), 0 reprise(s)).
\begin{corrigebox}[Corrigé extrait de la banque]
\CorrigeQuestion{1}
\begin{marginfigure}
\includegraphics[width=\linewidth]{14_cor_01}
\end{marginfigure}
\CorrigeQuestion{2}
Exprimons le couple du ressort par $C_r(\varphi)=a\varphi + b$. On a d'une part, $C_r(\varphi_0)=0$. D'autre part, on a une raideur $k$ de \SI{45}{Nm} pour un angle de rotation 100\degres soit $k=\dfrac{45}{100 \dfrac{\pi}{180}} = \SI{26}{Nm.rad^{-1}}$. On a donc  $C_r(\varphi_f)=k \dfrac{\pi}{2}$.

On a donc : 
$\left\{\begin{array}{l} 
a\varphi_0 + b = 0 \\
a\varphi_f + b = k \dfrac{\pi}{2}
\end{array}
\right. $
$\Leftrightarrow 
\left\{\begin{array}{l} 
 b = -a\varphi_0 \\
a\varphi_f  -a\varphi_0= k \dfrac{\pi}{2}
\end{array}
\right. $
$\Leftrightarrow 
\left\{\begin{array}{l} 
 b = -a\varphi_0 \\
a\left(\varphi_f  -\varphi_0\right)= k \dfrac{\pi}{2}
\end{array}
\right. $
$\Leftrightarrow 
\left\{\begin{array}{l} 
 b = -a\varphi_0 \\
a= k \dfrac{\pi}{2 \left(\varphi_f  -\varphi_0\right)}
\end{array}
\right. $

On a donc $C_r(\varphi)= k \dfrac{\pi}{2 \left(\varphi_f  -\varphi_0\right)}\varphi - k \dfrac{\pi \varphi_0}{2 \left(\varphi_f  -\varphi_0\right)} $.

Avec $\varphi_0=\dfrac{\pi}{2}$ et $\varphi_f=0$, on a 

$C_r(\varphi)= -k \varphi + k \dfrac{\pi}{2} $.
\CorrigeQuestion{3}
\begin{itemize}
\item On isole 1, on réalise un TMS en $C$ en projection sur $\vect{k_0}$. On obtient une équation liant le couple moteur et l'action normale dans la liaison sphère plan. 
\item On isole 2, on réalise un TMS en $A$ en projection sur $\vect{k_0}$. On obtient une équation liant le couple dans le ressort et l'action normale dans la liaison sphère plan. 
\item En combinant les deux équations on élimine l'action normale dans la liaison sphère plan. On peut éliminer un des deux angles en utilisant la loi entrée sortie.
\end{itemize}
\end{corrigebox}
